% Generated from Review/corrected_text.md (the reviewed source text).

\bakhsh{ب. ساختار مسابقات}

\band{۱. کلیات}

مسابقات در سه روز برگزار خواهد شد. اهداف دو روز اول مسابقات عبارت‌اند از:

\begin{itemize}
  \item انتخاب چهار بازیکنی که در روز سوم برای عنوان قهرمان جهان رقابت خواهند کرد.
  \item انتخاب چهار بازیکن زنی که در روز سوم برای عنوان قهرمان زنان جهان رقابت خواهند کرد (مگر آنکه نیازی نباشد؛ بند ج.۳ را ببینید).
  \item انتخاب چهار بازیکن نوجوانی که در روز سوم برای عنوان قهرمان نوجوانان جهان رقابت خواهند کرد (مگر آنکه نیازی نباشد؛ بند ج.۴ را ببینید).
  \item تعیین تیم قهرمان مسابقات جهانی اتللو.
\end{itemize}

روز سوم رقابت‌ها به شکل زیر برگزار خواهد شد:

\onvan{قهرمان جهان}

\begin{itemize}
  \item دو مسابقه نیمه‌نهایی برگزار خواهد شد: یکی بین نفرات اول و چهارم دو روز اول مسابقات و دیگری بین نفرات دوم و سوم. این مسابقات به شکل بهترین نتیجه از سه بازی خواهند بود.
  \item فینال بین برندگان دو مسابقه نیمه‌نهایی و به شکل بهترین نتیجه از سه بازی برگزار خواهد شد. همچنین یک بازی پلی‌آف بین بازندگان نیمه‌نهایی‌ها برای تعیین مقام سوم برگزار می‌شود.
\end{itemize}

\onvan{قهرمان زنان جهان}

\begin{itemize}
  \item دو مسابقه نیمه‌نهایی برگزار خواهد شد: یکی بین نفرات اول و چهارم زنان در دو روز اول مسابقات و دیگری بین نفرات دوم و سوم. این مسابقات به شکل تک‌بازی خواهند بود.
  \item فینال بین برندگان دو مسابقه نیمه‌نهایی زنان به شکل تک‌بازی برگزار خواهد شد و یک بازی پلی‌آف نیز بین بازندگان نیمه‌نهایی‌ها برای تعیین مقام سوم برگزار می‌شود.
\end{itemize}

\onvan{قهرمان نوجوانان جهان}

\begin{itemize}
  \item دو مسابقه نیمه‌نهایی برگزار خواهد شد: یکی بین نفرات اول و چهارم نوجوانان در دو روز اول مسابقات و دیگری بین نفرات دوم و سوم. این مسابقات به شکل تک‌بازی خواهند بود.
  \item فینال بین برندگان دو مسابقه نیمه‌نهایی نوجوانان به شکل تک‌بازی برگزار خواهد شد و یک بازی پلی‌آف نیز بین بازندگان نیمه‌نهایی‌ها برای تعیین مقام سوم برگزار می‌شود.
\end{itemize}

\band{۲. سیستم مسابقات}

در طول دو روز اول مسابقات، سیستم مورد استفاده به تعداد بازیکنان بستگی خواهد داشت:

\begin{itemize}
  \item تا ۸ بازیکن: دوره‌ای دو دور (\lr{Double Round Robin}).
  \item بین ۹ تا ۱۴ بازیکن: دوره‌ای یک دور (\lr{Single Round Robin}).
  \item بین ۱۵ تا ۲۴ بازیکن: سیستم سوئیسی اصلاح‌شده (\lr{Modified Swiss System}).
  \item بیش از ۲۴ بازیکن: سیستم سوئیسی اصلاح‌نشده (\lr{Unmodified Swiss System}).
\end{itemize}

برای اطلاعات بیشتر بند ب.۷ را ببینید. اگر از سیستم سوئیسی استفاده شود، هفت دور در روز اول و شش دور در روز دوم مسابقات برگزار خواهد شد. اگر تعداد بازیکنان فرد باشد، کشور میزبان حق دارد برای جلوگیری از استراحت (\lr{bye})، یک بازیکن اضافی وارد مسابقات کند. این بازیکن در ابتدای مسابقات به عنوان بازیکن اضافی آن کشور تعیین خواهد شد. چنانچه کشور میزبان به هر دلیلی نتواند تا پیش از آغاز مسابقات بازیکنی را معرفی کند، رئیس کل مسابقات می‌تواند یکی از بازیکنان غیرشرکت‌کننده را به عنوان بازیکن اضافی تعیین کند. این بازیکن در ابتدای مسابقات با عنوان «بازیکن اضافی \lr{WOF}» تعیین خواهد شد و برای قرعه‌کشی تحت عنوان «کشور \lr{WOF}» (\lr{WOF Country}) ثبت خواهد شد. امتیازات بازیکن اضافی هرگز برای قهرمانی تیمی محاسبه نخواهد شد. با این حال، این بازیکن در عنوان‌های انفرادی جایگاهی برابر با سایر بازیکنان خواهد داشت و به ویژه می‌تواند به مسابقات نیمه‌نهایی و فینال (و حسب مورد، مسابقات زنان و نوجوانان) راه یابد.

\band{۳. فهرست نهایی رده‌بندی فردی}

در پایان دو روز اول رقابت‌ها، مسئولان برگزاری فهرست نهایی رده‌بندی فردی را بر اساس روند زیر تهیه خواهند کرد:

\zirband{\LRE{۳\kern0pt.\kern0pt۱} فهرست موقت}

ابتدا یک فهرست موقت تهیه می‌شود که در آن بازیکنان بر اساس تعداد امتیازات خود در مسابقات رده‌بندی می‌شوند. برای هر برد یک امتیاز و برای هر تساوی نیم امتیاز داده می‌شود. در صورت تساوی امتیازات، برای تفکیک بازیکنان از سیستم ترکیبی برایتول (\lr{Mixed Brightwell System - MBS}) استفاده خواهد شد (پیوست الف را ببینید). اگر دو یا چند بازیکن در مسابقات امتیاز برابر داشته باشند، تفکیک آن‌ها در رده‌بندی به ترتیب نزولی ضریب برایتول (\lr{Brightwell Quotient - BQ}) انجام می‌شود. اگر دو یا چند بازیکن هم امتیاز برابر و هم ضریب برایتول برابر داشته باشند، تفکیک به ترتیب نزولی تعداد مهره‌ها انجام می‌شود. در صورت ادامه تساوی، تفکیک با یک روش تصادفی عادلانه انجام می‌شود (برای مثال پرتاب یک مهره اتللو یا بیرون کشیدن ژتون‌های شماره‌دار از یک کیسه دربسته).

\zirband{\LRE{۳\kern0pt.\kern0pt۲} پلی‌آف}

\begin{itemize}
  \item اگر بازیکنان رتبه چهارم و پنجم فهرست موقت امتیاز برابر داشته باشند، یک بازی پلی‌آف انجام خواهند داد.
  \item اگر بازیکنان رتبه سوم و ششم فهرست موقت امتیاز برابر داشته باشند، یک بازی پلی‌آف انجام خواهند داد.
  \item اگر بازیکنان رتبه دوم و هفتم فهرست موقت امتیاز برابر داشته باشند، یک بازی پلی‌آف انجام خواهند داد.
  \item اگر بازیکنان رتبه اول و هشتم فهرست موقت امتیاز برابر داشته باشند، یک بازی پلی‌آف انجام خواهند داد.
\end{itemize}

\onvan{انتخاب رنگ}

در هر بازی پلی‌آف، بازیکنی که رتبه بهتری دارد، حق انتخاب رنگ یا تساوی-برد را دارد (بند د.۹ را ببینید).

\zirband{\LRE{۳\kern0pt.\kern0pt۳} فهرست نهایی}

بازیکنانی که در روز سوم مسابقات در نیمه‌نهایی‌ها شرکت خواهند کرد عبارت‌اند از:

\begin{itemize}
  \item تمام بازیکنانی که در میان چهار نفر برتر فهرست موقت بودند و نیازی به انجام بازی پلی‌آف نداشتند.
  \item تمام برندگان بازی‌های پلی‌آف.
\end{itemize}

فهرست نهایی رده‌بندی فردی همان فهرست موقت است، با سه تغییر:

\begin{itemize}
  \item چهار بازیکنی که به نیمه‌نهایی راه یافته‌اند (بالا را ببینید)، به ترتیب جایگاهشان در فهرست موقت، رتبه‌های ۱ تا ۴ را می‌گیرند.
  \item تمام بازندگان بازی‌های پلی‌آف به طور مشترک در رتبه پنجم قرار می‌گیرند.
  \item بازیکنانی که در میان چهار نفر اول نیستند و امتیاز برابر دارند، هم‌رتبه اعلام می‌شوند. در صورت نیاز (برای مثال برای تقسیم جوایز نقدی)، این بازیکنان را می‌توان به ترتیب نزولی ضریب برایتول، به طور جداگانه رده‌بندی کرد (مانند فهرست رده‌بندی موقت).
\end{itemize}

\band{۴. فهرست نهایی رده‌بندی تیمی}

فهرست نهایی رده‌بندی برای تعیین تیم قهرمان جهان به این ترتیب تهیه می‌شود:

مسئولان برگزاری مجموع امتیازات بازیکنان هر تیم را محاسبه می‌کنند (در صورتی که تیمی چهار یا پنج بازیکن داشته باشد، فقط نتایج سه نفر برتر محاسبه می‌شود) و فهرست رده‌بندی بر اساس این مجموع تهیه خواهد شد (در صورتی که کشور میزبان یا \lr{WOF} یک بازیکن اضافی وارد مسابقات کرده باشد، امتیازات کسب‌شده توسط این بازیکن اضافی در مجموع تیمی در نظر گرفته نمی‌شود). در صورتی که دو یا چند تیم امتیاز برابر داشته باشند، تیمی در رده‌بندی بالاتر قرار می‌گیرد که بیشترین تعداد بازی را میان اعضای تیم خود انجام داده باشد (هیچ بازیکن اضافی عضو تیم محسوب نمی‌شود). اگر باز هم دو یا چند تیم برابر باشند، تیمی در اولویت قرار می‌گیرد که بهترین بازیکن را در فهرست نهایی رده‌بندی فردی دارد (بر اساس فهرست نهایی، که در آن برای تفکیک بازیکنان ابتدا از \lr{MBS} و سپس از تعداد مهره‌ها استفاده می‌شود). در صورت ادامه تساوی، تیمی در اولویت خواهد بود که بازیکن دوم آن در فهرست نهایی رده‌بندی فردی رتبه بالاتری دارد. در صورتی که هنوز نتوان این تیم‌ها را از هم تفکیک کرد، هم‌رتبه اعلام می‌شوند.

\band{۵. فهرست نهایی رده‌بندی زنان}

در پایان دو روز اول مسابقات، مسئولان برگزاری فهرست نهایی رده‌بندی زنان را بر اساس روند زیر تهیه می‌کنند:

\zirband{\LRE{۵\kern0pt.\kern0pt۱} فهرست موقت}

ابتدا یک فهرست موقت تهیه می‌شود که در آن بازیکنان زن بر اساس تعداد امتیازات خود در مسابقات رده‌بندی می‌شوند. برای هر برد یک امتیاز و برای هر تساوی نیم امتیاز داده می‌شود. در صورت تساوی امتیازات، برای تفکیک بازیکنان از سیستم ترکیبی برایتول (\lr{Mixed Brightwell System - MBS}) استفاده خواهد شد (پیوست الف را ببینید). اگر دو یا چند زن در مسابقات امتیاز برابر داشته باشند، تفکیک آن‌ها در رده‌بندی به ترتیب نزولی ضریب برایتول (\lr{BQ}) انجام می‌شود. اگر دو یا چند زن هم امتیاز برابر و هم ضریب برایتول برابر داشته باشند، تفکیک به ترتیب نزولی تعداد مهره‌ها انجام می‌شود. در صورت ادامه تساوی، تفکیک با یک روش تصادفی عادلانه انجام می‌شود (برای مثال پرتاب یک مهره اتللو یا بیرون کشیدن ژتون‌های شماره‌دار از یک کیسه دربسته).

\zirband{\LRE{۵\kern0pt.\kern0pt۲} فهرست نهایی}

اگر یک یا چند زن به نیمه‌نهایی مسابقات قهرمانی جهان راه یابند، نیمه‌نهایی زنان برگزار نخواهد شد. در عوض، تمام زنانی که در نیمه‌نهایی مسابقات قهرمانی جهان حضور دارند، بر اساس ترتیب نهایی خود پس از فینال مسابقات جهانی و پلی‌آف مقام سوم رده‌بندی خواهند شد. در غیر این صورت، زنانی که در روز سوم مسابقات در نیمه‌نهایی زنان شرکت خواهند کرد عبارت‌اند از:

\begin{itemize}
  \item تمام زنانی که در میان چهار نفر اول فهرست موقت هستند.
\end{itemize}

فهرست نهایی رده‌بندی زنان همان فهرست موقت است، با دو تغییر:

\begin{itemize}
  \item چهار زنی که به نیمه‌نهایی راه یافته‌اند، به ترتیب جایگاهشان در فهرست موقت، رتبه‌های ۱ تا ۴ را می‌گیرند.
  \item زنانی که در میان چهار نفر اول نیستند و امتیاز برابر دارند، هم‌رتبه اعلام می‌شوند.
\end{itemize}

در هر دو حالت، در صورت نیاز (برای مثال برای تقسیم جوایز نقدی)، زنانِ دارای امتیاز برابر در فهرست نهایی را می‌توان به طور جداگانه به ترتیب نزولی ضریب برایتول رده‌بندی کرد (مانند فهرست رده‌بندی موقت زنان).

\band{۶. فهرست نهایی رده‌بندی نوجوانان}

در پایان دو روز اول مسابقات، مسئولان برگزاری فهرست نهایی رده‌بندی نوجوانان را بر اساس روند زیر تهیه می‌کنند:

\zirband{\LRE{۶\kern0pt.\kern0pt۱} فهرست موقت}

ابتدا یک فهرست موقت تهیه می‌شود که در آن بازیکنان نوجوان بر اساس تعداد امتیازات خود در مسابقات رده‌بندی می‌شوند. برای هر برد یک امتیاز و برای هر تساوی نیم امتیاز داده می‌شود. در صورت تساوی امتیازات، برای تفکیک بازیکنان از سیستم ترکیبی برایتول (\lr{Mixed Brightwell System - MBS}) استفاده خواهد شد (پیوست الف را ببینید). اگر دو یا چند بازیکن نوجوان در مسابقات امتیاز برابر داشته باشند، تفکیک آن‌ها در رده‌بندی به ترتیب نزولی ضریب برایتول (\lr{BQ}) انجام می‌شود. اگر دو یا چند بازیکن نوجوان هم امتیاز برابر و هم ضریب برایتول برابر داشته باشند، تفکیک به ترتیب نزولی تعداد مهره‌ها انجام می‌شود. در صورت ادامه تساوی، تفکیک با یک روش تصادفی عادلانه انجام می‌شود (برای مثال پرتاب یک مهره اتللو یا بیرون کشیدن ژتون‌های شماره‌دار از یک کیسه دربسته).

\zirband{\LRE{۶\kern0pt.\kern0pt۲} فهرست نهایی}

اگر یک یا چند بازیکن نوجوان به نیمه‌نهایی مسابقات قهرمانی جهان راه یابند، نیمه‌نهایی نوجوانان برگزار نخواهد شد. در عوض، تمام بازیکنان نوجوانی که در نیمه‌نهایی مسابقات قهرمانی جهان حضور دارند، بر اساس ترتیب نهایی خود پس از فینال مسابقات جهانی و پلی‌آف مقام سوم رده‌بندی خواهند شد. در غیر این صورت، نوجوانانی که در روز سوم مسابقات در نیمه‌نهایی نوجوانان شرکت خواهند کرد عبارت‌اند از:

\begin{itemize}
  \item تمام بازیکنان نوجوانی که در میان چهار نفر اول فهرست موقت هستند.
\end{itemize}

فهرست نهایی رده‌بندی نوجوانان همان فهرست موقت است، با دو تغییر:

\begin{itemize}
  \item چهار بازیکن نوجوانی که به نیمه‌نهایی راه یافته‌اند، به ترتیب جایگاهشان در فهرست موقت، رتبه‌های ۱ تا ۴ را می‌گیرند.
  \item بازیکنان نوجوانی که در میان چهار نفر اول نیستند و امتیاز برابر دارند، هم‌رتبه اعلام می‌شوند.
\end{itemize}

در هر دو حالت، در صورت نیاز (برای مثال برای تقسیم جوایز نقدی)، بازیکنان نوجوانِ دارای امتیاز برابر در فهرست نهایی را می‌توان به طور جداگانه به ترتیب نزولی ضریب برایتول رده‌بندی کرد (مانند فهرست رده‌بندی موقت نوجوانان).

\band{۷. قرعه‌کشی در سیستم سوئیسی}

\zirband{\LRE{۷\kern0pt.\kern0pt۱} دور اول}

قرعه‌کشی دور اول به شرح زیر انجام می‌شود:

\begin{itemize}
  \item بازیکنان به دو گروه تقسیم می‌شوند: گروه \lr{A}، بازیکنانی که در ده دوره گذشته مسابقات جهانی اتللو دست‌کم در یکی از نیمه‌نهایی‌ها بازی کرده‌اند، و گروه \lr{B}، سایر بازیکنان.
  \item دو بازیکن به صورت تصادفی، یکی از گروه \lr{A} و دیگری از گروه \lr{B}، انتخاب می‌شوند. این دو بازیکن در دور اول با هم بازی خواهند کرد و بازیکن گروه \lr{A} با رنگ سیاه بازی می‌کند.
  \item این روند تا خالی شدن یکی از دو گروه تکرار می‌شود، با این تفاوت که رنگ بازیکن گروه \lr{A} به صورت یکی در میان سیاه و سفید خواهد بود.
  \item بازیکنان باقی‌مانده به صورت تصادفی قرعه‌کشی می‌شوند.
\end{itemize}

رئیس کل مسابقات مسئول روند دقیق قرعه‌کشی دور اول خواهد بود. او می‌تواند قرعه‌کشی را خودش انجام دهد یا از بازیکنان بخواهد که حریفان خود را به صورت تصادفی انتخاب کنند (یا از هر روش دیگری استفاده کند)، مشروط بر اینکه روند بالا رعایت شود.

\zirband{\LRE{۷\kern0pt.\kern0pt۲} دور دوم به بعد}

برای تعیین قرعه‌های دور دوم به بعد در سیستم سوئیسی، قوانین زیر باید به ترتیب اولویت (از بیشترین به کمترین) رعایت شوند:

\begin{bandha}
  \item[الف)] اگر از سیستم سوئیسی اصلاح‌نشده استفاده شود، دو بازیکنی که قبلاً یک بار با هم بازی کرده‌اند، دیگر نمی‌توانند با یکدیگر جفت شوند. اگر از سیستم سوئیسی اصلاح‌شده استفاده شود، دو بازیکنی که قبلاً دو بار با هم بازی کرده‌اند، دیگر نمی‌توانند با یکدیگر جفت شوند.
  \item[ب)] دو بازیکنی که با هم جفت می‌شوند باید امتیاز یکسانی داشته باشند؛ اگر این امر امکان‌پذیر نباشد، اختلاف امتیاز آن‌ها باید حداقل باشد.
  \item[ج)] در سیستم سوئیسی اصلاح‌شده، در صورت امکان، دو بازیکنی که قبلاً با یکدیگر بازی کرده‌اند، نباید دوباره با هم جفت شوند.
  \item[د)] بازیکنان باید در طول مسابقات به تعداد برابر با رنگ سیاه و سفید بازی کنند؛ اگر این امر امکان‌پذیر نباشد، اختلاف باید حداقل باشد.
  \item[هـ)] بازیکنی که قبلاً «بالا رفته» است (یعنی با حریفی از گروه امتیازی بالاتر جفت شده است)، باید «پایین بیاید» و بازیکنی که قبلاً «پایین رفته» است، باید «بالا برود».
  \item[و)] بازیکنان باید از هر دور به دور بعد رنگ خود را عوض کنند.
  \item[ز)] دو بازیکنی که برای یک کشور ثبت شده‌اند (تیم اصلی + بازیکن اضافی) نباید در سیستم سوئیسی با یکدیگر بازی کنند، به شرط آنکه بتوان بدون انجام این بازی، قرعه‌ای معادل (از نظر سایر قوانین بالا) یافت.
\end{bandha}

\band{۸. انصراف}

اگر بازیکنی از مسابقات انصراف دهد، دیگر نمی‌تواند به مسابقات بازگردد. به ویژه، او نمی‌تواند به هیچ نیمه‌نهایی یا فینالی راه یابد. با این حال، امتیاز او در زمان انصراف در فهرست نهایی رده‌بندی فردی ثبت می‌شود و در فهرست نهایی رده‌بندی تیمی لحاظ خواهد شد. بازیکنی که انصراف داده است، پس از انصراف در قرعه‌کشی‌ها در نظر گرفته نخواهد شد و به جای او یک بازیکن فرضی به نام \lr{BYE} وارد روند قرعه‌کشی خواهد شد. همچنین برای اثر بازیکنِ انصراف‌داده بر محاسبه \lr{MBS}، پیوست الف را ببینید.
